% ===== Preâmbulo compartilhado (Tectonic / XeTeX) — documentos da FTL094 =====
\documentclass[11pt,a4paper]{article}
\usepackage{fontspec}
\usepackage[brazil]{babel}
\usepackage{geometry}
\geometry{a4paper,margin=2.5cm}
\usepackage{amsmath,amssymb}
\usepackage{graphicx}
\usepackage{booktabs}
\usepackage{array}
\usepackage{longtable}
\usepackage{tabularx}
\usepackage{float}
\usepackage{enumitem}
\usepackage{xcolor}
\usepackage{tikz}
\usetikzlibrary{arrows.meta,positioning,automata,shapes.geometric,fit,backgrounds,calc}
\usepackage{pgfplots}
\pgfplotsset{compat=1.18}
\usepackage{listings}
\usepackage{caption}
\usepackage{fancyhdr}
\usepackage[hidelinks]{hyperref}

% ----- cores -----
\definecolor{ufamblue}{RGB}{0,51,102}
\definecolor{codegreen}{rgb}{0,0.5,0}
\definecolor{codeblue}{rgb}{0,0,0.7}
\definecolor{codered}{rgb}{0.6,0,0}
\definecolor{lgray}{rgb}{0.95,0.95,0.95}

% ----- listagens C -----
\lstdefinestyle{c}{language=C,basicstyle=\ttfamily\footnotesize,
  keywordstyle=\color{codeblue}\bfseries,commentstyle=\color{codegreen}\itshape,
  stringstyle=\color{codered},numbers=left,numberstyle=\tiny\color{gray},
  frame=single,breaklines=true,tabsize=2,showstringspaces=false,columns=fullflexible}

% ----- constantes do projeto -----
\newcommand{\projnome}{Sistema de Navegação Autônoma para o Robô Móvel Khepera~IV}
\newcommand{\disciplina}{FTL094 --- Engenharia de Software de Tempo Real}
\newcommand{\periodo}{Período Letivo 2026/1}
\newcommand{\versaodoc}{1.0}

% ----- cabeçalho/rodapé -----
\pagestyle{fancy}\fancyhf{}
\rhead{\small \tipodoc}
\lhead{\small Khepera~IV --- FTL094}
\cfoot{\thepage}
\renewcommand{\headrulewidth}{0.4pt}

% ----- coluna X centralizada p/ tabularx -----
\newcolumntype{C}{>{\centering\arraybackslash}X}

% ----- capa padronizada -----
\newcommand{\fazcapa}{%
\begin{titlepage}\centering
{\large\bfseries UNIVERSIDADE FEDERAL DO AMAZONAS (UFAM)}\\[3pt]
{\normalsize Centro de P\&D em Tecnologia Eletrônica e da Informação --- CETELI}\\[2pt]
{\normalsize Programa de Pós-Graduação em Engenharia Elétrica --- PPGEE}\\[26pt]
{\large \disciplina}\\[2pt]
{\normalsize \periodo}\\[64pt]
\rule{\textwidth}{1.2pt}\\[12pt]
{\Huge\bfseries\color{ufamblue} \tipodoc}\\[14pt]
{\Large \projnome}\\[10pt]
\rule{\textwidth}{1.2pt}\\[48pt]
{\large\bfseries Equipe}\\[8pt]
{\normalsize Gabriel Henrique de Souza Nazaré}\\[3pt]
{\normalsize Luan Alexandre Ramos Barreto}\\[3pt]
{\normalsize Matheus Rocha Canto}\\[3pt]
{\normalsize Matheus Serrão Uchôa}\\[3pt]
{\normalsize Ricardo Mendonça Braz}\\[3pt]
{\normalsize Yan Fernandes da Silva}\\[40pt]
{\normalsize Docente responsável: Prof. Dr.~[nome do docente]}\\
\vfill
{\small Documento de tipo \emph{\tipodoc} --- Versão \versaodoc}\\[2pt]
{\small Manaus --- AM, \datadoc}
\end{titlepage}}

% ----- tabela de controle de versões -----
\newcommand{\controleversoes}{%
\section*{Controle de versões}
\begin{tabularx}{\textwidth}{@{}l l l X@{}}
\toprule
\textbf{Versão} & \textbf{Data} & \textbf{Autor} & \textbf{Descrição} \\
\midrule
1.0 & \datadoc & Equipe FTL094 & Emissão inicial do documento. \\
\bottomrule
\end{tabularx}
\bigskip}

\newcommand{\tipodoc}{Plano de Testes}
\newcommand{\datadoc}{30 de junho de 2026}

\begin{document}
\fazcapa
\controleversoes
\tableofcontents
\newpage

% =====================================================================
\section{Introdução}
% =====================================================================
\subsection{Objetivo}
Este documento define a \textbf{estratégia, os casos e os critérios de teste} do \emph{\projnome}, de modo a verificar, de forma sistemática e rastreável, o atendimento aos requisitos especificados. É o artefato de planejamento que orienta a posterior execução, cujos resultados são registrados no documento de \emph{Protocolo de Testes}.

\subsection{Escopo}
Abrange testes \textbf{funcionais}, \textbf{não funcionais} (segurança/desempenho) e de \textbf{tempo real} do software de controle, conduzidos no simulador Webots. Não abrange, nesta fase, a validação no robô físico nem ambientes com obstáculos móveis.

\subsection{Referências}
Documento de Análise de Requisitos; Documento de Arquitetura; I.~Sommerville, \emph{Engenharia de Software} (cap.~8); IEEE Std 829 (\emph{Test Documentation}).

% =====================================================================
\section{Estratégia de testes}
% =====================================================================
\subsection{Níveis de teste}
\begin{description}[leftmargin=2.2cm,style=nextline]
  \item[Unidade] Verificação das funções auxiliares isoladas (normalização de ângulo \texttt{wrap\_pi}, saturação \texttt{clampd}, cálculo de erro de rumo e cinemática diferencial), por inspeção e asserções pontuais.
  \item[Integração] Verificação do encadeamento Percepção\,$\to$\,FSM\,$\to$\,Controle\,$\to$\,Atuação dentro de um ciclo do laço.
  \item[Sistema] Execução da \emph{missão completa} A\,$\to$\,B no cenário com obstáculos.
\end{description}

\subsection{Abordagem}
Predomina o teste \textbf{caixa-preta orientado a requisitos} (verifica-se o comportamento observável frente a cada requisito), complementado por \textbf{caixa-branca} (inspeção do laço de controle e das transições da máquina de estados). A execução é \textbf{automatizável em modo \emph{headless}} e a observação se dá por \textbf{telemetria instrumentada} (arquivo \texttt{run.log}), o que confere objetividade e reprodutibilidade.

\subsection{Procedimento padrão de execução}
\begin{enumerate}[noitemsep,topsep=2pt]
  \item configurar o mundo (arena, obstáculos, posições inicial e alvo);
  \item compilar o controlador e executar o Webots em modo \texttt{--batch --mode=fast --no-rendering};
  \item coletar a telemetria periódica em \texttt{run.log};
  \item comparar o resultado observado com o esperado e registrar PASSOU/FALHOU.
\end{enumerate}

% =====================================================================
\section{Itens a testar}
% =====================================================================
\begin{itemize}[noitemsep,topsep=2pt]
  \item \textbf{A testar:} RF-01 a RF-07; RNF-01 (segurança); RT-01 (período do laço).
  \item \textbf{A verificar por inspeção:} RF-08 (configuração) e RF-09 (sinalização de estado).
  \item \textbf{Fora desta fase:} portabilidade ao robô físico (RNF-04) e obstáculos dinâmicos.
\end{itemize}

% =====================================================================
\section{Ambiente de teste}
% =====================================================================
\begin{table}[H]\centering
\caption{Configuração do ambiente de teste.}
\begin{tabularx}{\textwidth}{@{}l X@{}}
\toprule
\textbf{Item} & \textbf{Configuração} \\
\midrule
Simulador & Webots R2025a (modo \emph{headless}: \texttt{--batch --mode=fast --no-rendering}). \\
Controlador & \texttt{goto\_avoid} (linguagem C, FSM Navegar/Girar/Contornar). \\
Cenário & Arena $3\times3$~m; 5 obstáculos (\texttt{WoodenBox}); início $(-1,-1)$; alvo $(1,1)$. \\
Instrumentação & Telemetria periódica em \texttt{run.log} (descarga por linha via \texttt{fflush}). \\
Localização & GPS + IMU (simulação). \\
\bottomrule
\end{tabularx}
\end{table}

% =====================================================================
\section{Critérios gerais de aprovação/reprovação}
% =====================================================================
\textbf{Aprovado (PASSOU)} quando o resultado observado corresponde ao esperado \emph{e} não há colisão. \textbf{Reprovado (FALHOU)} em caso de colisão, não-chegada ao alvo dentro do tempo previsto, ou violação do período de tempo real. Defeitos identificados são registrados e, após correção, o caso é \emph{reexecutado}.

% =====================================================================
\section{Casos de teste}
% =====================================================================
\begin{table}[H]\centering
\caption{Especificação dos casos de teste.}
\renewcommand{\arraystretch}{1.25}
\begin{tabularx}{\textwidth}{@{}l X l X@{}}
\toprule
\textbf{ID} & \textbf{Objetivo} & \textbf{Req.} & \textbf{Resultado esperado} \\
\midrule
CT-01 & Navegar em caminho livre (sem obstáculos no trajeto direto). & RF-01, RF-03 & Robô atinge o alvo e para dentro da tolerância (7~cm). \\
CT-02 & Desviar de um obstáculo frontal isolado. & RF-02 & Robô contorna o obstáculo e retoma o curso ao alvo. \\
CT-03 & Concluir a missão no \emph{slalom} de 5 obstáculos. & RF-02, RNF-01 & Robô chega ao alvo \textbf{sem colisões}. \\
CT-04 & Escapar de mínimo local (alvo atrás de obstáculo frontal). & RF-06 & Robô não trava; contorna e prossegue. \\
CT-05 & Respeitar o período do laço de controle. & RT-01 & Laço periódico de 16~ms sem perda de passo. \\
CT-06 & Registrar telemetria e calibração dos sensores. & RF-07 & \texttt{run.log} gerado; base de IR caracterizada por medição. \\
\bottomrule
\end{tabularx}
\end{table}

\noindent\textbf{Detalhamento --- CT-03 (caso principal).}
\emph{Pré-condições:} mundo com 5 obstáculos; robô em $(-1,-1)$ orientado ao alvo. \emph{Passos:} executar a missão em modo \emph{headless}; coletar \texttt{run.log}. \emph{Entrada:} alvo $(1,1)$. \emph{Resultado esperado:} distância ao alvo decresce até $<7$~cm; nenhuma leitura indica colisão; sequência de estados \textsc{Navegar}/\textsc{Girar}/\textsc{Contornar} observada. \emph{Critério:} chegada sem colisão.

\noindent\textbf{Detalhamento --- CT-04 (robustez).}
\emph{Pré-condições:} obstáculo posicionado entre robô e alvo (rumo $\approx 0$). \emph{Resultado esperado:} o robô \emph{não} permanece preso (posição não estaciona); a manobra de escape é acionada e a missão prossegue.

% =====================================================================
\section{Critérios de suspensão e retomada}
% =====================================================================
Os testes são \textbf{suspensos} diante de defeito bloqueante (impede a continuidade da missão) ou indisponibilidade do ambiente. São \textbf{retomados} após a correção do defeito, com \emph{reexecução} dos casos afetados e de regressão dos demais.

% =====================================================================
\section{Responsáveis e cronograma}
% =====================================================================
A execução é liderada pelo Engenheiro de Testes (Luan Alexandre Ramos Barreto), com apoio da equipe, e deve concluir-se até \textbf{30/06/2026}, com os resultados consolidados no Protocolo de Testes.

% =====================================================================
\section{Riscos de teste}
% =====================================================================
\begin{itemize}[noitemsep,topsep=2pt]
  \item \textbf{Cobertura geométrica limitada:} nem toda configuração de obstáculos é exercitada (mitigação: incluir cenário de mínimo local, CT-04).
  \item \textbf{Divergência simulação\,$\to$\,real:} resultados em simulação podem não refletir o robô físico (mitigação: registrar como ressalva; teste físico em fase futura).
  \item \textbf{Ausência de medição formal de WCET:} o período é verificado por regularidade da telemetria, não por análise de pior caso (mitigação: indicar como trabalho futuro).
\end{itemize}

\end{document}
